\begin{description}
\item[(1)] The student should be able to import the ascii data of the
  \emph{txt} file in the \emph{MS Excel} spreadsheet application.
  \hfill(12 Points)
\item[(2)] In order to calculate the angular distance $\varphi$ (next step),
  all coordinates should be converted into decimal degrees. The students
  should be able to convert the [hours, min, sec] and the
  [$^\circ$,\ $'$,\ $''$] into decimal degrees, using. % sic
  For the right ascension:
  $\alpha_{\deg} = h\times15 + \dfrac{m}{60} + \dfrac{s}{3600}$
  and for the declination:
  $\delta_{\deg} = {}^\circ + \dfrac{'}{60} + \dfrac{''}{3600}$.
  All coordinates are positive, so the student does not have to worry about
  checking the sign, which is not trivial. \hfill(12 Points)
\item[(3)] Using the cosine law,
  $\varphi = \arccos(\sin\delta_1\times\sin\delta_2
    + \cos\delta_1\times\cos\delta_2\times\cos\,(a_1 - a_2))$,
  the angular distance, $\varphi$, between each of the stars and the point of
  convergence should be calculated \hfill(15 Points)
\item[(4)] The student should easily calculate the proper motion of each
  star, by inserting the given equation,
  $\mu = \sqrt{(\mu_a\cos\delta)^2 + \mu_d^2}$, in the spreadsheet.
  \hfill(9 Points)
\item[(5)] Again, inserting the given equation,
  $r_\mu = \dfrac{v_r\tan\varphi}{4.74047\,\mu}$ in the spreadsheet
  (remembering to divide the given values of $\mu$ by 1000 to get arcseconds),
  the student should be able to calculate the distance, $r_\mu$, of each star.
  Any star whose distance from the centre of the cluster is larger than
  10\,pc should be omitted from the following calculations. \hfill(12 Points)
\item[(6)] The trigonometric parallax distance is given by
  $r_\pi = \dfrac{1}{\pi\times10^{-3}}$. This equation should be inserted in
  the spreadsheet. The result is in parsecs [pc]. \hfill(6 Points)
\item[(7)] Using the statistical functions \emph{AVERAGE} and \emph{STDEV} of
  MS Excel, the student should easily calculate the \emph{average} and the
  \emph{standard deviation} of $r_\mu$ and $r_\pi$. \hfill(6 Points)
\item[(8)] The method with the smaller standard deviation is, obviously more
  accurate. \hfill(3 Point)
\end{description}

(\emph{macros} are allowed)
